\section{Excluding models with several directions}\label{sec:bounds}

We now combine the logarithmic inequality with estimates on smooth
approximants of a model.  All curvature quantities in this section belong
to the approximants; no differentiability of the limiting body is used.

\begin{proposition}\label{prop:one-direction}
Suppose \(3\le n\le9\).  The largest number of directions of a model in
\(\mathcal F(D)\), in the sense of Definition~\ref{def:model}, is one.
\end{proposition}

We first establish the area estimate that prevents the measures in the
logarithmic inequality from vanishing. The comparison uses concavity
of the affine-area integrand and stationarity, as in the local
maximization argument of Trudinger and Wang
\cite[Section~6, Equations~(6.2)--(6.4)]{TW2000}. Here the comparison
is applied to uniformly normalized caps of smooth approximants.

\begin{lemma}\label{lem:positive-cap-area}
Let \(D_j\) be affine images of \(D\) converging locally to a
full-dimensional closed convex set \(C\).  Suppose that a nonzero linear
function \(\ell\) and a number \(b\) satisfy
\[
 C\cap\{\ell\le b\}\ \text{is compact},
 \qquad
 \operatorname{int}C\cap\{\ell<b\}\ne\varnothing.
\]
There are \(c>0\) and \(j_0\) such that
\[
 \int_{\partial D_j\cap\{\ell<b\}}\dd A_j\ge c
 \qquad(j\ge j_0).
\]
\end{lemma}

\begin{proof}
Lemma~\ref{lem:cap-compactness} gives a uniform bound for these caps.
Choose a closed ball \(\overline B_r(o)\) in
\(\operatorname{int}C\cap\{\ell<b\}\).
After decreasing \(r\), the same ball lies in \(D_j\) for all sufficiently
large \(j\).  Fix \(\varepsilon_0>0\) with
\(\ell(o)\le b-\varepsilon_0\).

Let \(e_j\) be the unit vector in the image of the original upward vertical
direction under the affine map defining \(D_j\).  It is a recession
direction of \(D_j\), and \(\partial D_j\) is graphical along it.  The
component \(\ell(e_j)\) is bounded below by a positive constant.  Indeed,
the ray \(o+t e_j\) lies in \(D_j\); if \(\ell(e_j)\le0\), the whole ray
lies in the cap, and otherwise its initial segment of length
\((b-\ell(o))/\ell(e_j)\) lies there.  Uniform boundedness of the caps
excludes both a nonpositive component and components tending to zero.
Consequently
\[
 V_j=\frac{e_j}{\ell(e_j)}
 \quad\text{satisfies}\quad
 \ell(V_j)=1,\qquad |V_j|\le C.
\]
Use the coordinates
\[
 X=\zeta+\tau V_j,\qquad \zeta\in\ker\ell.
\]
These coordinate maps and their inverses are uniformly bounded.
With an orthonormal basis of \(\ker\ell\), their determinants have absolute
value \(1/|\ell|\), where \(|\ell|\) is the Euclidean norm of the covector.
In these coordinates \(\partial D_j\) is the graph of a smooth strictly
convex function \(g_j\), and the cap corresponds to \(\tau\le b\).

The sets
\[
 E_j=\{\zeta:g_j(\zeta)<b\}
\]
are uniformly bounded.  Every closed sublevel of \(g_j\) at a level at
most \(b\) is compact within its domain: its lifted graph lies in a
compact cap, and the closed hypersurface \(\partial D_j\) is precisely
the graph in this direction.
Writing \(\zeta_j=o-\ell(o)V_j\), the common interior ball gives
\[
 g_j(\zeta)\le b-\varepsilon_0
 \qquad\bigl(\zeta\in B_{\ker\ell}(\zeta_j,r)\bigr).
\]
Here \(B_{\ker\ell}(\zeta_j,r)\) is the radius-\(r\) ball in
\(\ker\ell\) centered at \(\zeta_j\).
The centers \(\zeta_j\) and the sets \(E_j\) are uniformly bounded.
We may therefore choose one \(\alpha>0\) so small that the quadratics
\[
 q_j(\zeta)=b-\frac{2\varepsilon_0}{3}
                +\alpha|\zeta-\zeta_j|^2
\]
satisfy
\[
 b-\frac{2\varepsilon_0}{3}\le q_j
       \le b-\frac{\varepsilon_0}{3}
 \qquad\text{on }E_j.
\]

Let \(\chi\) be a smooth convex regularization of the positive part, with
\(0\le\chi'\le1\), equal to \(0\) for arguments at most \(-\eta\) and
equal to the argument for arguments at least \(\eta\), where
\(0<\eta<\varepsilon_0/6\).  On \(E_j\) set
\[
 \widetilde g_j=g_j+\chi(q_j-g_j).
\]
This equals \(q_j\) on the indicated base ball and equals \(g_j\) near
the level \(b\).  Extending by \(g_j\) outside \(E_j\) gives a smooth
perturbation whose difference from \(g_j\) has compact support in \(E_j\).
Its Hessian is positive definite, since it is a convex combination of
\(D^2g_j\) and \(D^2q_j\), plus a positive semidefinite matrix.

The function \(F(H)=(\det H)^\delta\) is concave on positive definite
matrices: the determinant \(n\)-th root is concave, and
\(n\delta=n/(n+2)<1\).
Concavity and Lemma~\ref{lem:stationarity} give
\begin{align*}
 \int_{E_j}F(D^2\widetilde g_j)\,\dd\zeta
 &\le \int_{E_j}F(D^2g_j)\,\dd\zeta\\
 &\quad+\delta\int_{E_j}
 F(D^2g_j)(D^2g_j)^{-1}:D^2(\widetilde g_j-g_j)\,\dd\zeta\\
 &=\int_{E_j}F(D^2g_j)\,\dd\zeta.
\end{align*}
The first integral is bounded below by
\(\lvert B_{\ker\ell}(0,r)\rvert(2\alpha)^{n\delta}>0\).
Affine covariance of area, with the uniformly controlled coordinate
determinants above, proves the assertion.
\end{proof}

\subsection{Positive mass away from the coordinate faces}

Fix a model \(C\) with the maximal number \(k\) of directions, and choose
affine images \(D_j\to C\).  Write its coordinates as
\((s,y)\in\R^k\times\R^m\), where \(m=n+1-k\).
Choose \(b>0\) so that the cap for
\(\ell_0(s)=s_1+\cdots+s_k\) contains an interior ball strictly below its
top.  Such a choice is possible because \(C\) has nonempty interior.
The cap is compact by Definition~\ref{def:model}.
Lemma~\ref{lem:positive-cap-area} gives
\begin{equation}\label{eq:cap-area-positive}
 A_j\bigl(\partial D_j\cap\{\ell_0<b\}\bigr)\ge c>0
\end{equation}
for large \(j\), and all these caps lie in a fixed ball.

For any fixed \(0<\varepsilon<1\), convergence and the uniform cap bound
imply that \(s_i\ge-\varepsilon\) everywhere on the approximant caps
once \(j\) is sufficiently large.  In the part where \(s_i\le\varepsilon\),
stretch the single coordinate \(s_i\) by \(1/\varepsilon\).
The transformed boundary part lies in a fixed ball, independent of
\(\varepsilon\) and \(j\).  Lemma~\ref{lem:area-bound} and the
\(|\det L|^{n\delta}\) covariance of affine area give
\begin{equation}\label{eq:strip-area}
 A_j\bigl(\partial D_j\cap\{\ell_0<b,\ s_i\le\varepsilon\}\bigr)
       \le C\varepsilon^{n\delta}.
\end{equation}
Only the boundary part under consideration needs to lie in the ball.
Choose \(\varepsilon\) so small that the sum of these \(k\) bounds is
less than \(c/2\), and then choose \(j\) large enough.
Equations~\eqref{eq:cap-area-positive} and~\eqref{eq:strip-area} leave
area at least \(c/2\) in a fixed bounded region where all \(s_i>\varepsilon\).

If \(m=0\), the model is the nonnegative orthant in \(\R^{n+1}\).
The remaining region is contained in a compact subset of its interior,
which lies in \(\operatorname{int}D_j\) for large \(j\).
It cannot contain any points of \(\partial D_j\), a contradiction.
Hence
\begin{equation}\label{eq:transverse-dimension}
 k\le n,\qquad m\ge1.
\end{equation}

For the rest of the section, suppose that \(k\ge2\).
On every fixed compact box of positive \(s\)'s, the fibers of \(D_j\)
are uniformly bounded and contain zero in their interiors for large
\(j\).  This follows from Proposition~\ref{prop:centering},
Lemma~\ref{lem:model-fibers}, and Lemma~\ref{lem:cap-compactness}.
Thus the tube coordinates and measures
\eqref{eq:tube-area}, \eqref{eq:tube-volume}, and \eqref{eq:log-measure}
are available there.
For a base set \(E\subset(0,\infty)^k\), write \(\mu_j(E)\) and
\(\sigma_j(E)\) for the masses of \(E\times S^{m-1}\).
In particular, define
\[
 Q_L=\{s\in(0,\infty)^k:|\log s|_\infty\le L\},
 \qquad
 \log s=(\log s_1,\ldots,\log s_k).
\]

The preceding area bound gives a fixed positive box \(Q\) such that
\(\int_{Q\times S^{m-1}}I_0\,\dd s\,\dd\omega\ge c/2\) for large \(j\).
On this box \(h\) is uniformly bounded above, and
\[
 \int_{Q\times S^{m-1}}\det R\,\dd s\,\dd\omega\le C.
\]
Indeed, \(\det R\,\dd\omega\) is Euclidean area on the boundary of a
uniformly bounded convex fiber.  When \(m=1\), it is counting measure
on the two endpoints, with \(\det R=1\).
Using the density \(\mu_0\) from \eqref{eq:tube-volume}, the exact
factorization
\[
 I_0=\mu_0^{2\delta}h^{n\delta}(\det R)^{1-2\delta}
\]
and H\"older's inequality yield
\[
 \frac c2\le\int_{Q\times S^{m-1}} I_0\,\dd s\,\dd\omega
       \le C\mu_j(Q)^{2\delta}.
\]
Since \(\sigma_j\ge\mu_j\), enlarging \(Q\) to a fixed \(Q_{L_0}\)
proves that there are \(L_0,c_0>0\) with
\begin{equation}\label{eq:positive-log-mass}
 \sigma_j(Q_{L_0})\ge c_0
 \qquad\text{for all sufficiently large }j.
\end{equation}

\subsection{A uniform bound on each logarithmic cell}

We now have positive mass in one fixed logarithmic box. To control
growth, we need a bound for the mass of every cell that is independent
of its logarithmic position. Centering provides this uniformity after
affine normalization. The cap identity controls an inverse-Hessian
integral, and bounds for Hessian minors provide the complementary
factor in a weighted interpolation.

\begin{lemma}\label{lem:log-cell-bound}
There is a constant \(C_{\mathrm{cell}}\), independent of
\(a\in\R^k\), such that
\[
 \sigma_j\left\{s:\tfrac12e^{a_i}\le s_i\le2e^{a_i}
                       \text{ for every }i\right\}
 \le C_{\mathrm{cell}}
\]
for every sufficiently large \(j\).  The threshold for \(j\) may depend
on \(a\).
\end{lemma}

\begin{proof}
Rescale \(s_i\) by \(e^{-a_i}\) and make an invertible linear change in
\(y\).  By Corollary~\ref{cor:normalized-fibers}, these changes can be chosen
so that the limit fiber above \(\mathbf1=(1,\ldots,1)\) satisfies
\[
 B_{r_*}\subset K_{\mathbf1}\subset B_{R_*},
\]
where \(r_*,R_*>0\) are independent of \(a\).
We use the same affine change on the approximants.
The measures \(\sigma_j\) are preserved by these normalizations.

Corollary~\ref{cor:normalized-fibers} also gives uniform fiber bounds on
the fixed box \([1/4,4]^k\), uniform bounds for caps with fixed positive
linear forms and fixed top levels, and a fixed interior ball centered at
\(o=(\mathbf1,0)\).  Applying its approximation statement on a slightly
larger positive box gives
\begin{equation}\label{eq:normalized-fiber-bounds}
 0<c_1\le h\le C_1
 \quad\text{on }[1/4,4]^k\times S^{m-1}
\end{equation}
and the same uniform cap and interior-ball bounds for all sufficiently
large \(j\).  These constants are independent of \(a\); only the
required index depends on the normalization.

Let \(Q=[1/2,2]^k\) and \(\mathcal T=Q\times S^{m-1}\).
By \eqref{eq:log-measure} and \eqref{eq:normalized-fiber-bounds},
\[
 \sigma_j(Q)\le C\int_{\mathcal T}
       \sqrt{\det B\,\det R}\,(1+\tr B^{-1})
       \,\dd s\,\dd\omega,
\]
since \(s_i^{-2}\le4\) on \(Q\).  To control the inverse-Hessian
factor in this integral, we first prove
\begin{equation}\label{eq:cap-inverse-hessian}
 \int_{\mathcal T} I_0\,\tr B^{-1}\,\dd s\,\dd\omega\le C.
\end{equation}
For \(r=1,\ldots,k\), set
\[
 a^{(r)}=\mathbf1+e_r,\qquad \ell_r(s)=a^{(r)}\cdot s.
\]
Their values on \(Q\) lie in
\(J=[(k+1)/2,\,2(k+1)]\).
Choose a smooth nonnegative function \(\rho\), compactly supported
below a fixed number \(b_0>2(k+1)\), with \(\rho\ge1\) on \(J\),
and put
\[
 \psi(t)=\int_t^{b_0}(r-t)\rho(r)\,\dd r.
\]
Then \(\psi''=\rho\ge0\), and \(\psi\) vanishes above \(b_0\).
Although it is affine sufficiently far below, its surface test is
compactly supported in the bounded cap \(\ell_r\le b_0\).
The functions \(\psi,\psi'\) and the levels \(\ell_r\) are bounded there
by uniform constants.

Apply Lemma~\ref{lem:cap-identity} with the center \(o\) of the fixed
interior ball.  If its radius is \(r_0\), the supporting-plane inequality
gives
\[
 \nu(o-X)\ge r_0|\nu|.
\]
The integrand involving \(\psi''\) in that identity is nonnegative over
the whole cap.  On the tube, take the conormal
\(\nu=(D_sh,-e)\); it has norm at least one, and
\[
 |d\ell_r|_{H^{-1}}^2=(a^{(r)})^{\mathsf T}B^{-1}a^{(r)}.
\]
The other terms in the cap identity have bounded integrals by
Lemma~\ref{lem:area-bound} and the uniform cap bound.  Hence
\[
 \int_{\mathcal T} I_0\,
       (a^{(r)})^{\mathsf T}B^{-1}a^{(r)}
       \,\dd s\,\dd\omega\le C.
\]
Summing over \(r\) proves \eqref{eq:cap-inverse-hessian}, since
\[
 \sum_{r=1}^k a^{(r)}(a^{(r)})^{\mathsf T}
       =\Id+(k+2)\mathbf1\mathbf1^{\mathsf T}\ \ge\ \Id.
\]

We next claim that
\begin{equation}\label{eq:hessian-minor-bound}
 \int_{\mathcal T}\det B\,(1+\tr B^{-1})\,\dd s\,\dd\omega\le C.
\end{equation}
For each fixed \(e\), the function \(-h(\,\cdot\,,e)\) is strictly convex.
Its gradient is uniformly bounded on \(Q\) by
\eqref{eq:normalized-fiber-bounds} on the larger box.
For a principal minor, hold the unused \(s\)-coordinates fixed and
consider the gradient in the selected coordinates.  It is injective,
because the corresponding Hessian is positive definite, and its image
lies in a fixed bounded box.  The change-of-variables formula bounds the
integral of its Jacobian, which is precisely that principal minor.
Integrating in the unused coordinates and in \(e\) gives a uniform bound.
The determinant itself is the full principal minor, and
\(\det B\,\tr B^{-1}\) is the sum of the principal minors of order \(k-1\).
This proves \eqref{eq:hessian-minor-bound}.
As before, the uniform fiber bound also gives
\begin{equation}\label{eq:fiber-area-bound}
 \int_{\mathcal T}\det R\,\dd s\,\dd\omega\le C.
\end{equation}

Equations~\eqref{eq:hessian-minor-bound} and
\eqref{eq:fiber-area-bound}, together with Cauchy--Schwarz, give
\[
 \int_{\mathcal T}\sqrt{\det B\,\det R}
       \,\dd s\,\dd\omega\le C.
\]
For the inverse-Hessian weight, set
\[
 \theta=\frac1{2(1-\delta)}=\frac{n+2}{2(n+1)},
 \qquad t_B=\tr B^{-1}.
\]
The interpolation identity is
\begin{equation}\label{eq:weighted-interpolation}
 \sqrt{\det B\,\det R}\,t_B
       =(I_0t_B)^\theta(\det B\,t_B)^{1-\theta}.
\end{equation}
Since \(0<\theta<1\), H\"older's inequality applied to
\eqref{eq:weighted-interpolation}, using
\eqref{eq:cap-inverse-hessian} and~\eqref{eq:hessian-minor-bound},
gives
\[
 \int_{\mathcal T}\sqrt{\det B\,\det R}\,
           \tr B^{-1}\,\dd s\,\dd\omega\le C.
\]
The last two bounds control both terms in the target integral and hence
\(\sigma_j(Q)\).
Its invariance under the preceding coordinate changes proves the lemma.
\end{proof}

\subsection{The growth contradiction}

\begin{proof}[Proof of Proposition~\ref{prop:one-direction}]
A model with one direction exists by Lemma~\ref{lem:proper-sections}.
Suppose that the maximal number \(k\) is at least two, and fix the model
and approximants used above.  The case \(k=n+1\) has already been
excluded in \eqref{eq:transverse-dimension}; hence \(2\le k\le n\).

Covering a logarithmic box by finitely many of the cells in
Lemma~\ref{lem:log-cell-bound} gives a constant \(C_{\mathrm{box}}\)
such that, for every fixed \(L\ge1\),
\begin{equation}\label{eq:polynomial-log-growth}
 \sigma_j(Q_L)\le C_{\mathrm{box}}(L+1)^k
 \qquad\text{for every sufficiently large }j.
\end{equation}
The constant is independent of \(L\) and \(j\).  The index threshold may
depend on \(L\), because only finitely many cells are involved for each
fixed \(L\).  By increasing that threshold, the tube coordinates are
also available on \(Q_{L+2}\).

For an integer \(l\ge1\), choose a smooth cutoff
\(\widetilde v_l\) on \(\R^k\), equal to one on
\(\{|t|_\infty\le l\}\), zero outside
\(\{|t|_\infty<l+1\}\), and with a uniformly bounded Euclidean gradient.
Its gradient is supported in the intervening shell.
Proposition~\ref{prop:log-inequality}, applied with
\(v_l(s)=\widetilde v_l(\log s)\), implies
\[
 \sigma_j(Q_l)
       \le C_{\mathrm{grow}}
            \bigl(\sigma_j(Q_{l+1})-\sigma_j(Q_l)\bigr)
\]
whenever the tube is available through \(Q_{l+1}\).
Here \(C_{\mathrm{grow}}\) is independent of \(l\) and \(j\).
Box boundaries have zero mass, because these measures have smooth
densities in the tube coordinates.
Writing \(q=1+C_{\mathrm{grow}}^{-1}>1\) and taking an integer
\(l_0\ge\max(1,L_0)\), iteration from
\eqref{eq:positive-log-mass} gives
\begin{equation}\label{eq:exponential-log-growth}
 \sigma_j(Q_L)\ge c_0 q^{L-l_0}
 \qquad(L\ge l_0,\ L\text{ an integer})
\end{equation}
for approximants on which all the indicated estimates hold.

Choose one finite integer \(L\) large enough that
\[
 c_0q^{L-l_0}>C_{\mathrm{box}}(L+1)^k.
\]
Then choose one \(j\) large enough for the initial mass bound, the
finitely many cell bounds through \(Q_L\), and tube coordinates through
\(Q_{L+2}\).  For this same \(j\),
\eqref{eq:polynomial-log-growth} and~\eqref{eq:exponential-log-growth}
contradict each other.  Thus no maximal model has \(k\ge2\), proving
the proposition.
\end{proof}
